% =============================================================================
% 4.4 ARQUITECTURA DE COMUNICACIÓN EN TIEMPO REAL
% =============================================================================
\section{Arquitectura de Comunicación en Tiempo Real}
\label{sec:realtime}

La plataforma ofrece chat y notificaciones \textbf{en tiempo real} sin implementar un
servidor de \textit{WebSockets} dentro de Spring. Esta decisión —documentada como
ADR (sección~\ref{sec:adr}, pág.~\pageref{sec:adr})— delega la entrega \textit{push} en
\textbf{Supabase Realtime}, conservando la autoridad de validación en el backend.

\subsection{Infraestructura de mensajería en tiempo real}
\label{subsec:rt_infra}

El patrón es de \textbf{escritura validada por el backend y entrega \textit{push} por el
canal Realtime}. Combina lo mejor de dos mundos: la autoridad del servidor (validación,
idempotencia, persistencia) y la latencia mínima de la entrega directa al cliente
suscrito. La figura~\ref{fig:rt-arq} ilustra la topología.

\vspace{0.3cm}
\begin{figure}[h!]
\centering
\begin{tikzpicture}[node distance=1.0cm and 1.8cm,
    box/.style={rectangle, rounded corners=3pt, draw, font=\sffamily\scriptsize, align=center,
        minimum width=2.7cm, minimum height=0.95cm, inner sep=3pt},
    rel/.style={-{Stealth[length=2mm]}, semithick}]
    \node[box, fill=c4Persona!25] (emisor) {Emisor\\\scriptsize (SPA)};
    \node[box, fill=c4Componente!45, right=of emisor] (be) {ChatController\\\scriptsize valida + persiste};
    \node[box, fill=c4Datos!22, draw=c4Datos!70!black, right=of be] (db) {PostgreSQL\\\scriptsize core.chat\_messages};
    \node[box, fill=colorAcento!28, draw=colorAcento!80!black, below=of db] (rt) {Supabase\\Realtime};
    \node[box, fill=c4Persona!25, below=of emisor] (dest) {Destinatario\\\scriptsize (SPA suscrita)};

    \draw[rel] (emisor) -- node[c4etiqueta, above]{\scriptsize POST /messages} (be);
    \draw[rel] (be) -- node[c4etiqueta, above]{\scriptsize INSERT} (db);
    \draw[rel] (db) -- node[c4etiqueta, right]{\scriptsize WAL / trigger} (rt);
    \draw[rel] (rt) -- node[c4etiqueta, above]{\scriptsize push (sin polling)} (dest);
\end{tikzpicture}
\caption{Topología de mensajería: escritura validada por el backend, entrega \textit{push} por Supabase Realtime.}
\label{fig:rt-arq}
\end{figure}

El flujo concreto de un mensaje es:

\begin{enumerate}[noitemsep]
    \item El emisor envía el mensaje por REST a \comando{POST /api/v1/messages}.
    \item \comando{ChatController} valida que el llamante sea \textbf{participante} de la
    conversación y persiste el mensaje con \textbf{idempotencia} (clave de cliente,
    migración V7), evitando duplicados ante reintentos de red.
    \item \textbf{Supabase Realtime} detecta el cambio en la tabla mediante el registro
    de escritura anticipada (WAL) y lo entrega al destinatario suscrito, \textbf{sin
    \textit{polling}}.
\end{enumerate}

\begin{NotaTecnica}
La idempotencia es esencial en redes móviles inestables: si el cliente reintenta un
\texttt{POST} cuya respuesta se perdió, el backend reconoce la clave de idempotencia y
devuelve el mensaje ya persistido en lugar de crear un duplicado. Esta garantía se
implementa con una restricción \texttt{UNIQUE} parcial en la base de datos
(véase el capítulo~\ref{ch:basedatos}, pág.~\pageref{ch:basedatos}).
\end{NotaTecnica}

\subsection{Máquina de estados de una sesión de chat}
\label{subsec:rt_estados}

El ciclo de vida de un mensaje puede modelarse como una máquina de estados finitos. La
figura~\ref{fig:fsm-chat} la representa, incluyendo el \textbf{borrado lógico} que
preserva la integridad histórica (los mensajes no se eliminan físicamente, solo se marcan
con \texttt{deleted\_at}).

\vspace{0.3cm}
\begin{figure}[h!]
\centering
\begin{tikzpicture}[node distance=1.4cm and 2.0cm,
    estado/.style={rectangle, rounded corners=8pt, draw=c4Sistema!70!black, fill=c4Componente!30,
        font=\sffamily\scriptsize\bfseries, align=center, minimum width=2.2cm, minimum height=0.8cm},
    ini/.style={circle, fill=flujoInicio, minimum size=0.4cm, inner sep=0pt},
    rel/.style={-{Stealth[length=2mm]}, semithick}]
    \node[ini] (start) {};
    \node[estado, right=1.2cm of start] (compose) {Redactando};
    \node[estado, right=of compose] (sent) {Enviado\\(persistido)};
    \node[estado, right=of sent] (delivered) {Entregado\\(Realtime)};
    \node[estado, below=of delivered] (read) {Leído};
    \node[estado, below=of sent] (deleted) {Borrado lógico};

    \draw[rel] (start) -- (compose);
    \draw[rel] (compose) -- node[c4etiqueta, above]{POST /messages} (sent);
    \draw[rel] (sent) -- node[c4etiqueta, above]{push} (delivered);
    \draw[rel] (delivered) -- node[c4etiqueta, right]{visto} (read);
    \draw[rel] (sent) -- node[c4etiqueta, left]{eliminar} (deleted);
    \draw[rel] (delivered) -- node[c4etiqueta, above, sloped]{eliminar} (deleted);
\end{tikzpicture}
\caption{Máquina de estados de un mensaje de chat, con borrado lógico.}
\label{fig:fsm-chat}
\end{figure}

Las \textbf{notificaciones} siguen un patrón análogo y asíncrono: se crean en respuesta a
eventos (mensaje nuevo, \textit{like} recibido, etc.) y el usuario las consume vía
\comando{GET /api/v1/notifications}, marcándolas como leídas individualmente
(\comando{PATCH /\{id\}/read}) o en bloque (\comando{POST /read-all}). La tabla
\texttt{notifications} dispara, a su vez, eventos Realtime que actualizan el contador del
frontend sin recargar la página.

\subsection{Notificaciones por correo asíncronas (SMTP)}
\label{subsec:email}

Más allá de las notificaciones in-app, el backend envía \textbf{correos transaccionales}
mediante \comando{EmailService}, que usa el cliente de Spring Mail sobre SMTP
(\texttt{smtp.gmail.com:587} con STARTTLS). Los correos cubren los eventos de la
tabla~\ref{tab:emails}. El envío se realiza de forma que no bloquee la respuesta HTTP al
usuario, preservando la latencia percibida.

\vspace{0.2cm}
\noindent
\renewcommand{\arraystretch}{1.3}
\fontsize{8.5}{10.2}\selectfont\sffamily
\begin{tabularx}{\textwidth}{|L{3.6cm}|Y|}
\hline
\rowcolor{headerTabla}
\sffamily\textbf{Evento} & \sffamily\textbf{Correo enviado} \\
\hline
Recuperación de contraseña & Código OTP de un solo uso (sección~\ref{subsec:seq_otp}, pág.~\pageref{subsec:seq_otp}). \\ \hline
\rowcolor{grisTecnico}
Cambio de correo & Confirmación con enlace a \pathbreak{/api/v1/auth/confirm-email-change}. \\ \hline
Eventos administrativos & Comunicaciones del panel de administración (\pathbreak{/api/admin/email}). \\ \hline
\end{tabularx}
\label{tab:emails}

\subsection{Preocupaciones transversales del servidor}
\label{subsec:transversales}

Tres preocupaciones atraviesan todos los dominios del backend y se resuelven de forma
centralizada, no por dominio:

\begin{itemize}[noitemsep]
    \item \textbf{Idempotencia}: las operaciones de escritura sensibles a reintentos (envío
    de mensajes) aceptan una clave de idempotencia que evita duplicados, respaldada por una
    restricción de unicidad en la base de datos.
    \item \textbf{Registro estructurado} (\textit{logging}): cada servicio emite trazas con
    SLF4J; los errores no recuperables se registran con \textit{stack trace} completo en el
    servidor, mientras que al cliente solo viaja un mensaje seguro.
    \item \textbf{Zona horaria y serialización}: Jackson se configura globalmente en UTC y
    con \texttt{default-property-inclusion: NON\_NULL}, garantizando respuestas compactas y
    marcas de tiempo coherentes en toda la API.
\end{itemize}
